\begin{trapbox}
	\begin{description}[style=nextline, leftmargin=2.8em, labelsep=0.5em, itemsep=0.35em]
		\item[\textbf{\textcolor{CTrap}{易错点一（比例颠倒）}}] 错误地认为 $\frac{BD}{DC}=\frac{AC}{AB}$ 或 $\frac{BD}{DC}=\frac{AB}{BC}$。
		\item[\textbf{\textcolor{CTrap}{正确理解（比例对应）：}}] 分子 $BD$ 对应边 $AB$，分母 $DC$ 对应边 $AC$，即 $\frac{BD}{DC}=\frac{AB}{AC}$。
		\item[\textbf{\textcolor{CTrap}{错因分析（记忆失真）：}}] 死记结论，未从面积法推导中理解对应关系，导致分子分母混淆。
	\end{description}

	\medskip
	\begin{description}[style=nextline, leftmargin=2.8em, labelsep=0.5em, itemsep=0.35em]
		\item[\textbf{\textcolor{CTrap}{易错点二（条件误用）}}] 没有确认 $AD$ 是否为角平分线，仅凭图形相似或直观就套用定理。
		\item[\textbf{\textcolor{CTrap}{正确理解（前提检查）：}}] 使用定理前必须明确 $AD$ 平分 $\angle BAC$，或通过已知条件推出。
		\item[\textbf{\textcolor{CTrap}{错因分析（审题粗心）：}}] 忽略关键条件，将一般分点误当作角平分线。
	\end{description}

	\medskip
	\begin{description}[style=nextline, leftmargin=2.8em, labelsep=0.5em, itemsep=0.35em]
		\item[\textbf{\textcolor{CTrap}{易错点三（计算失误）}}] 解 $\frac{x}{a-x}=\frac{m}{n}$ 时，交叉相乘出现符号错误或漏乘。
		\item[\textbf{\textcolor{CTrap}{正确理解（乘积方程）：}}] 应化为 $n x = m (a-x)$，移项注意变号。
		\item[\textbf{\textcolor{CTrap}{错因分析（运算不牢）：}}] 交叉相乘后括号展开或移项时符号处理错误。
	\end{description}
\end{trapbox}
